\subsection{Area of a Circle}
%
%
\renewcommand{\theequation}{\theenumi}
\begin{enumerate}[label=\arabic*.,ref=\thesubsection.\theenumi]
\numberwithin{equation}{enumi}
%
%
%\item
%	Using Fig. \ref{ch5_sin_theta}, show that 
%	%
%\begin{equation}
%\label{ch5_sin_theta_eq}
%\sin  \theta_1 = \sin \brak{\theta_1 + \theta_2}\cos \theta_2 - \cos\brak{\theta_1+\theta_2}\sin\theta_2
%\end{equation}	
%	%
%
%\begin{figure}[!ht]
%	\begin{center}
%		
%		%\includegraphics[width=\columnwidth]{./figs/ch5_sin_theta}
%		%\vspace*{-10cm}
%		\resizebox{\columnwidth}{!}{\begin{tikzpicture}
[scale =3,>=stealth,point/.style = {draw, circle, fill = black, inner sep = 1pt},]

\node (A) at (0,3)[point,label=above :$A$] {};
\node (B) at (3,0)[point,label=below :$B$] {};
\node (C) at (0,0)[point,label=below :$C$] {};
\node (D) at (0,1.5)[point,label=left :$D$] {};
\draw (A)--(B);
\draw (C)--(B);
\draw (A)--(C);
\draw (B)--(D);
\tkzMarkAngle[size=.4](A,B,D);
\tkzMarkAngle[size=.3](D,B,C);
\tkzMarkRightAngle[size=.15](A,C,B);

\node [above] at (1.6,1.5){$c$};
\node [below] at (1.6,0){$a$};
\node [above] at (-0.2,1.5){$b$};
\node [above] at (2.5,0){$\theta_2$};
\node [above] at (2.5,0.3){$\theta_1$};
\end{tikzpicture}}
%	\end{center}
%	\caption{$\sin2\theta = 2\sin\theta\cos\theta$}
%	\label{ch5_sin_theta}	
%\end{figure}
%%
%
%\solution The following equations can be obtained from the figure using the forumula for the area of a triangle
%%
%\begin{align}
%ar \brak{\Delta ABC} &= \frac{1}{2}ac \sin\brak{\theta_1 + \theta_2} \\
%&= ar \brak{\Delta BDC} + ar \brak{\Delta ADB} \\
%&= \frac{1}{2}cl \sin{\theta_1} + \frac{1}{2}al \sin{\theta_2} \\ 
%&= \frac{1}{2}ac \sin{\theta_1} \sec \theta_2 + \frac{1}{2}a^2 \tan{\theta_2} 
%\end{align}
%$\brak{\because
%	l = a \sec \theta_2}$.  From the above,
%\begin{align}
%\Rightarrow \sin\brak{\theta_1 + \theta_2} &=  \sin{\theta_1} \sec \theta_2 + \frac{a}{c} \tan{\theta_2} \\
%\Rightarrow \sin\brak{\theta_1 + \theta_2} &=  \sin{\theta_1} \sec \theta_2 + \cos\brak{\theta_1 + \theta_2} \tan{\theta_2} 
%\end{align}
%Multiplying both sides by $\cos \theta_2$,
%\begin{align}
%\Rightarrow \sin\brak{\theta_1 + \theta_2}\cos{\theta_2} &=  \sin{\theta_1}  + \cos\brak{\theta_1 + \theta_2} \sin\theta_2  
%\end{align}
%%
%resulting in
%\begin{equation}
%\Rightarrow \sin \theta_1 = \sin\brak{\theta_1 + \theta_2}\cos{\theta_2} - \cos\brak{\theta_1 + \theta_2} \sin\theta_2 
%\end{equation}
%\item
%	Prove the following identities 
%	%
%	\begin{enumerate}
%\item 
%\begin{equation}
%		\label{ch5_sin_diff}
%\sin\brak{\alpha - \beta} = \sin \alpha \cos \beta - \cos \alpha \sin \beta.
%\end{equation}
%\item 
%\begin{equation}
%\cos\brak{\alpha + \beta} = \cos \alpha \cos \beta - \sin \alpha \sin \beta.
%		\label{ch5_cos_diff}
%\end{equation}
%
%	\end{enumerate}
%	%
%
%\solution In \eqref{ch5_sin_theta_eq}, let
%%
%\begin{equation}
%\begin{split}
%\theta_1 + \theta_2 &= \alpha \\
%\theta_2 &=  \beta
%\end{split}
%\end{equation}
%%
%This gives \eqref{ch5_sin_diff}.  In \eqref{ch5_sin_diff}, replace $\alpha$ by 
%%
%$90{\degree} - \alpha$.  This results in
%%
%\begin{multline}
%\sin\brak{90{\degree} - \alpha - \beta}
%\\
%=
%\sin \brak{90{\degree} -\alpha} \cos \beta - \cos \brak{90{\degree} -\alpha} \sin \beta \\
%\Rightarrow \cos\brak{\alpha + \beta} = \cos \alpha \cos \beta - \sin \alpha \sin \beta
%\end{multline}
%% 
%\item
%	Using \eqref{ch5_sin_theta_eq} and \eqref{ch5_cos_diff}, show that
%\begin{align}
%\label{ch5_sin_sum}
%\sin\brak{\theta_1 + \theta_2} &= \sin\theta_1  \cos\theta_2 + \cos\theta_1\sin\theta_2
%\\
%\cos\brak{\theta_1 - \theta_2} &= \cos\theta_1  \cos\theta_2  \sin\theta_1\sin\theta_2
%\label{ch5_cos_sum}
%\end{align}
%
%%
%\solution From \eqref{ch5_sin_theta_eq},
%%
%\begin{align}
% \sin \brak{\theta_1 + \theta_2}\cos \theta_2 =\sin  \theta_1 +\cos\brak{\theta_1+\theta_2}\sin\theta_2 
%\end{align}
%%
%Using \eqref{ch5_cos_diff} in the above,
%%
%\begin{multline}
%\sin \brak{\theta_1 + \theta_2}\cos \theta_2 
%=\sin  \theta_1 +\lbrak{\cos \theta_1\cos\theta_2 }
%\\	
%\rbrak{	- \sin \theta_1\sin\theta_2}\sin\theta_2 
%\end{multline}
%%
%which can be expressed as
%%
%\begin{multline}
%\sin \brak{\theta_1 + \theta_2}\cos \theta_2 
%=\sin  \theta_1 +\cos \theta_1\cos\theta_2 \sin\theta_2 
%\\	
%	- \sin \theta_1\sin^2\theta_2
%\end{multline}
%%
%Since
%%
%\begin{equation}
%\sin^2\theta_2 = 1- \cos^2\theta_2, 
%\end{equation}
%%
%we obtain
%%
%\begin{multline}
%\sin \brak{\theta_1 + \theta_2}\cos \theta_2 
%=\cos \theta_1\cos\theta_2 \sin\theta_2 
%\\	
%+ \sin \theta_1\cos^2\theta_2
%\end{multline}
%%
%resulting in
%%
%\begin{equation}
%\sin \brak{\theta_1 + \theta_2}
%=\cos \theta_1 \sin\theta_2 
%+ \sin \theta_1\cos\theta_2
%\end{equation}
%%
%after factoring out $\cos \theta_2$.  Using a similar approach, \eqref{ch5_cos_sum} can also be proved.
%%
%\item
%	Show that
%	%
%	\begin{equation}
%	\label{eq:sin2theta}
%	\sin 2\theta = 2 \sin\theta \cos\theta
%	\end{equation}
%	%
%%
%\item
%	Show that
%	%
%	\begin{align}
%	\label{eq:cos2theta_cos_sin}
%	\cos 2\theta &= \cos^2\theta -\sin^2\theta 
%\\
%	\label{eq:cos2theta_sin_sq}
%&= 1 - \sin^2\theta
%\\
%&=2\cos^2\theta -1
%	\label{eq:cos2theta_cos_sq}
%	\end{align}
%	%
\item The ratio of the perimeter of a circle to its diameter is $\pi$.
\label{prob:circ_peri_dia}
\item {\em Radian} is a another unit of the angle defined by
%
\begin{align}
\pi \text{ radians} = 180 \degree
\end{align}
%
\item
	In Fig. \ref{ch5_polygon_def}, 6 congruent triangles are arranged in a circular fashion.  Such a figure is known as a regular hexagon.  In general, $n$ number of traingles can be arranged to form a regular polygon.
\begin{figure}[!ht]
	\begin{center}
		
		%\includegraphics[width=\columnwidth]{./figs/ch5_polygon_def}
		%\vspace*{-10cm}
		\resizebox{\columnwidth}{!}{\begin{tikzpicture}
[scale =2,>=stealth,point/.style = {draw, circle, fill = black, inner sep = 1pt},]

\node (O) at (0,0)[point,label=above :$O$] {};
\node (F) at (-3,0)[point,label=left :$F$] {};
\node (C) at (3,0)[point,label=right :$C$] {};
\node (E) at (-1.5,3)[point,label=above :$E$] {};
\node (D) at (1.5,3)[point,label=above :$D$] {};
\node (A) at (-1.5,-3)[point,label=below :$A$] {};
\node (B) at (1.5,-3)[point,label=below :$B$] {};
\draw (A)--(D);
\draw (B)--(E);
\draw (F)--(C);
\draw (A)--(F);
\draw (F)--(E);
\draw (E)--(D);
\draw (D)--(C);
\draw (C)--(B);
\draw (B)--(A);

\tkzMarkAngle[size=.2](F,O,A);
\tkzMarkAngle[size=.22](E,O,F);
\tkzMarkAngle[size=.26](D,O,E);
\tkzMarkAngle[size=.2](C,O,D);
\tkzMarkAngle[size=.22](B,O,C);
\tkzMarkAngle[size=.26](A,O,B);

\node [above] at (-0.8,-1.5){$r$};
\node [above] at (-1.5,0){$r$};
\node [above] at (-0.9,1.5){$r$};
\node [above] at (0.9,1.5){$r$};
\node [above] at (1.5,0){$r$};
\node [above] at (0.8,-1.5){$r$};
\node [above] at (0,-0.7){$\frac{2\pi}{6}$};

\end{tikzpicture}
}
	\end{center}
	\caption{Polygon Definition}
	\label{ch5_polygon_def}	
\end{figure}
%
\item
The angle formed by each of the congruent triangles at the centre of a regular polygon of $n$ sides is $\frac{2\pi}{n} = \frac{2\pi}{n}$ rad.
%
\item 	The triangle that forms a polygon of $n$ sides is given in Fig. \ref{ch5_polygon_area}. Show that 
%
\begin{equation}
BC = 2r \sin\frac{\pi}{n}
\label{eq:circ_chord_len}
\end{equation}
%

\begin{figure}[!ht]
	\begin{center}
		
		%\includegraphics[width=\columnwidth]{./figs/ch5_polygon_area}
		%\vspace*{-10cm}
		\resizebox{\columnwidth}{!}{\begin{tikzpicture}
[scale =2,>=stealth,point/.style = {draw, circle, fill = black, inner sep = 1pt},]

\node (A) at (0,3)[point,label=above :$A$] {};
\node (B) at (-3,0)[point,label=below :$B$] {};
\node (C) at (3,0)[point,label=below :$C$] {};
\draw (A)--(B);
\draw (C)--(B);
\draw (A)--(C);

\tkzMarkAngle[size=.2](B,A,C);
\node [above] at (0.05,2.5){$\frac{2\pi}{n}$};
\node [above] at (-1.6,1.5){$r$};
\node [above] at (1.6,1.5){$r$};
\end{tikzpicture}
}
	\end{center}
	\caption{Triangle that forms a polygon}
	\label{ch5_polygon_area}	
\end{figure}
%
\solution From  \eqref{eq:tri_crad_R}.
%
\begin{equation}
BC = 2r \sin \frac{A}{2} = 2r\sin\frac{\pi}{n}
\end{equation}
%%
\item
Show that the perimeter of a regular polygon is given by 
%
\begin{equation}
\label{eq:peri_poly_n}
2rn \sin\frac{\pi}{n}
\end{equation}
%
\item
Show that the area of a regular polygon is given by 
%
\begin{equation}
\frac{n}{2}r^{2}\sin\frac{2\pi}{n}
\end{equation}
%
\solution  From Fig. 	\ref{ch5_polygon_area}	

%
\begin{equation}
\begin{split}
ar\brak{polygon} &= n \times ar\brak{\Delta ABC} \\
&= \frac{n}{2}r^{2}\sin\frac{2\pi}{n}
\end{split}
\end{equation}
%
using \eqref{eq:circ_area_chord}

\item
	Using Fig. \ref{ch5_circle_squeeze}, show that
%
\begin{equation}
\label{ch5_circle_squeeze_eq}
\frac{n}{2}r^{2}\sin\frac{2\pi}{n} < \text{ area of circle } < nr^{2}\tan\frac{\pi}{n}
\end{equation}
%
The portion of the circle visible in Fig. \ref{ch5_circle_squeeze} is defined to be a sector of the circle.

\begin{figure}[!ht]
	\begin{center}
		
		%\includegraphics[width=\columnwidth]{./figs/ch5_circle_squeeze}
		%\vspace*{-10cm}
		\resizebox{\columnwidth}{!}{\begin{tikzpicture}
[scale =2,>=stealth,point/.style = {draw, circle, fill = black, inner sep = 1pt},]

\def\rad{3}
\coordinate [point, label={right: $O$ }] (O) at (0, 0);
\draw (O) circle (\rad);
\node (P) at (0,-3)[point,label=below :$P$] {};
\node (A) at (-2,-2.22)[point,label= left :$A$] {};
\node (B) at (2,-2.22)[point,label=right :$B$] {};
\node (C) at (-2.7,-3)[point,label=below :$C$] {};
\node (D) at (2.7,-3)[point,label=below :$D$] {};
\draw (A)--(B);
\draw (O)--(B);
\draw (A)--(O);
\draw (O)--(P);
\draw (C)--(P);
\draw (D)--(P);

\draw [thick,dashed](A)--(C);
\draw [thick,dashed](B)--(D);
\tkzMarkRightAngle[size=0.2](O,P,D)
\tkzMarkAngle[size=0.2](C,O,P)
\tkzMarkAngle[size=0.3](P,O,D)
\node [above] at (-1.4,-1.5){$r$};
\node [above right] at (0,-1.5){$r$};
\node [above] at (1.4,-1.5){$r$};
\node [above right] at (0.4,-0.5){$\theta=360/n$};
\end{tikzpicture}}
	\end{center}
	\caption{Circle Area in between Area of Two Polygons}
	\label{ch5_circle_squeeze}	
\end{figure}
%

\solution Note that the circle is squeezed between the inner and outer regular polygons.  As we can see from Fig. \ref{ch5_circle_squeeze}, the area of the circle should be in between the areas of the inner and outer polygons.  Since
%
\begin{align}
ar \brak{\Delta OAB} &= \frac{1}{2}r^{2}\sin\frac{2\pi}{n} \\
ar \brak{\Delta OPQ} &= 2 \times \frac{1}{2} \times r \tan\frac{2\pi/n}{2} \times r \\
&= r^{2}\tan\frac{\pi}{n},
\end{align}
%
we obtain \eqref{ch5_circle_squeeze_eq}.
%
\item
Show that
	%
	\begin{equation}
	\label{ch5_circle_squeeze_simple}
\cos^2\frac{\pi}{n} < \frac{\text{ area of circle }}{nr^{2}\tan\frac{\pi}{n}} < 1	\end{equation}
	%

\solution From \eqref{ch5_circle_squeeze_eq} and \eqref{eq:sin2theta},
	%
{\small
	\begin{align}
	\frac{n}{2}r^{2}\sin\frac{2\pi}{n} < \text{ area of circle } 
%	\\
	< nr^{2}\tan\frac{\pi}{n} 
	\\
\Rightarrow 	
	{n}r^{2}\sin\frac{\pi}{n}\cos\frac{\pi}{n} < \text{ area of circle } 
%	\\
	< nr^{2}\tan\frac{\pi}{n} 
	\end{align}
%
}
which yields 	\eqref{ch5_circle_squeeze_simple} upon making use of the fact that 
%
\begin{align}
\frac{\sin \theta}{\cos \theta} = \tan \theta
\end{align}
%

%
\item
	Show that 
	\begin{equation}
	\label{ch5_cos_zero}
	\cos 0^{\degree} = 1
	\end{equation}

\solution Follows from the fact that $\cos 0 \degree = \sin \brak{90\degree -0\degree} = \sin \brak{90\degree }=1$ using \eqref{eq:sin90} and \eqref{eq:tri_90-ang}.



\item
	Show that 
	\begin{equation}
	\label{ch5_sin_zero}
	\sin 0^{\degree} = 0
	\end{equation}
%\solution Follows from the fact that $\sin 0 = 0$ and \eqref{eq:tri_sin_cos_id}.

\item
	Show that for large values of $n$
	%
	\begin{equation}
\label{eq:cos_zero_lim}
	%
\cos^2\frac{\pi}{n} = 1
%
	\end{equation}	
	% 

%
\solution  As $n \to \infty, \frac{\pi}{n} \to 0$. From \eqref{ch5_cos_zero}, this yields \eqref{eq:cos_zero_lim}.

%
\item  \eqref{eq:cos_zero_lim} is a {\em limit} and 
	 expressed as 
%
\begin{equation}
\label{eq:cos_zero_lim_def}
\lim_{n \rightarrow \infty}\cos^2\frac{\pi}{n} = 1
\end{equation}
%	

\item
	Show that 
	%
\begin{equation}
\label{ch5_circle_squeeze_tan_n}
\text{ area of circle } = r^2\lim_{n \rightarrow \infty}
{n\tan\frac{\pi}{n}} 
	%
\end{equation}	
	% 
\solution From \eqref{ch5_circle_squeeze_simple} and \eqref{eq:cos_zero_lim_def}, 
	%
	\begin{align}
%	\label{ch5_circle_squeeze_simple}
\lim_{n\to \infty}\cos^2\frac{\pi}{n} < \lim_{n\to \infty} \frac{\text{ area of circle }}{nr^{2}\tan\frac{\pi}{n}} < 1	
\\
1 = \lim_{n\to \infty} \frac{\text{ area of circle }}{nr^{2}\tan\frac{\pi}{n}} < 1	
\end{align}
	%
resulting in \eqref{ch5_circle_squeeze_tan_n}.
%
\item Show that 

	\begin{align}
\label{eq:peri_poly_sin_n}
 \lim_{n\to \infty}n \sin\frac{\pi}{n} &= \pi 
\end{align}
%
\solution From \eqref{prob:circ_peri_dia} and \eqref{eq:peri_poly_n}, the perimeter of the circle is 
%
\begin{align}
%\label{eq:peri_poly_sin_n}
\lim_{n\to \infty}2rn \sin\frac{\pi}{n} &= 2\pi r
\implies \lim_{n\to \infty}n \sin\frac{\pi}{n} &= \pi 
\end{align} 
\item Show that 

	\begin{equation}
\label{eq:peri_poly_tan_n}
	\pi = \lim_{n \rightarrow \infty}
	{n\tan\frac{\pi}{n}}
	\end{equation}
\solution 
%
From Fig. \eqref{ch5_circle_squeeze}, using the fact that the inner and outer polygons converge into a circle for large $n$,
\begin{align}
%\label{eq:peri_poly_n}
\lim_{n\to \infty} nCD -nAB &= 0
\\
\implies \lim_{n\to \infty} 2r n\tan\frac{\pi}{n}-2r n\sin\frac{\pi}{n} &= 0
\end{align}
%
from which, we obtain \eqref{eq:peri_poly_tan_n}
by substituting from \eqref{eq:peri_poly_sin_n}.


\item Show that the area of a circle is $\pi r^2$.
\\
\solution Use \eqref{eq:peri_poly_tan_n} in \eqref{ch5_circle_squeeze_tan_n}.

%\item
%	The radian is a unit of angle defined by
%\begin{equation}
%	1 \text{ radian} = \frac{2\pi}{2\pi}
%\end{equation}

%
%\item
%	Show that the circumference of a circle is $2 \pi r$.
\item
	Show that
	\begin{equation}
	\lim_{\theta \rightarrow 0} \frac{\sin\theta}{\theta} =
	\lim_{\theta \rightarrow 0} \frac{\tan\theta}{\theta} = 1
	\end{equation}

\item
	Show that the area of a sector with angle $\theta$ in radians is $\frac{1}{2}r^2\theta$.


\end{enumerate}
